\documentclass[aspectratio=169,xcolor=dvipsnames]{beamer}
\usetheme{Berlin}

\usepackage[english]{babel}
\usepackage{hyperref}
\usepackage{graphicx}
\usepackage{booktabs}
\usepackage{amsmath}
\usepackage{setspace}
\setbeamertemplate{caption}[numbered]
\usepackage[dvipsnames,svgnames,x11names]{xcolor}
\usepackage{xurl}
\usepackage{algorithm}
\usepackage{algorithmicx}
\usepackage{algpseudocode}
\usepackage{adjustbox}

\hypersetup{
    colorlinks=true,
    linkcolor=cyan,
    filecolor=blue,
    urlcolor=blue,
    citecolor=blue,
}

%----------------------------------------------------------------------------------------
%   CODE LISTINGS SETTINGS
%----------------------------------------------------------------------------------------
\usepackage{listings}
\definecolor{backcolour}{rgb}{0.97,0.97,0.99}
\definecolor{codegreen}{rgb}{0,0.6,0}

\lstdefinestyle{MATLABStyle}{
  language=Matlab,
  basicstyle=\ttfamily\scriptsize,
  keywordstyle=\color{blue}\bfseries,
  commentstyle=\color{codegreen},
  stringstyle=\color{violet},
  breaklines=true,
  numbers=left,
  numbersep=5pt,
  frame=lines,
  backgroundcolor=\color{backcolour}
}

\lstdefinestyle{PythonStyle}{
  language=Python,
  basicstyle=\ttfamily\scriptsize,
  keywordstyle=\color{blue}\bfseries,
  commentstyle=\color{codegreen},
  stringstyle=\color{red},
  breaklines=true,
  numbers=left,
  numbersep=5pt,
  frame=lines,
  backgroundcolor=\color{backcolour}
}

\lstset{style=MATLABStyle}

%----------------------------------------------------------------------------------------
%   TITLE
%----------------------------------------------------------------------------------------
\title{Ciclos \texttt{while}}
\subtitle{Algoritmos y Programaci\'on}
\author{Prof. D.Sc. BARSEKH-ONJI Aboud}
\institute{
    Facultad de Ingenier\'ia \\
    Universidad An\'ahuac M\'exico
}
\date{\today}

%----------------------------------------------------------------------------------------
%   AUTO AGENDA
%----------------------------------------------------------------------------------------
\AtBeginSection[]{
  \begin{frame}{Agenda}
    \tableofcontents[currentsection]
  \end{frame}
}

\begin{document}

\begin{frame}
    \titlepage
\end{frame}

%------------------------------------------------
\begin{frame}{Contenido general}
    \tableofcontents
\end{frame}

%================================================
\section{Introducci\'on a los ciclos}
%================================================

\begin{frame}{?`Por qu\'e necesitamos ciclos?}
    \begin{block}{Motivaci\'on}
        En programaci\'on, muchas tareas requieren \textbf{repetir} un conjunto de instrucciones
        varias veces. Sin ciclos, tendr\'iamos que escribir el mismo c\'odigo una y otra vez.
    \end{block}
    \vspace{0.3cm}
    \textbf{Ejemplos de la vida real:}
    \begin{itemize}
        \item Revisar todos los correos de una bandeja de entrada.
        \item Calcular la suma de 1000 n\'umeros.
        \item Pedir al usuario que ingrese datos hasta que sean v\'alidos.
        \item Actualizar una simulaci\'on en cada paso de tiempo.
    \end{itemize}
\end{frame}

\begin{frame}{Tipos de ciclos}
    \begin{columns}
        \column{0.5\textwidth}
        \textbf{Ciclo \texttt{for}}
        \begin{itemize}
            \item Se usa cuando se conoce \textbf{de antemano} el n\'umero de iteraciones.
            \item Ejemplo: recorrer una lista, iterar $n$ veces.
        \end{itemize}
        \vspace{0.3cm}
        \begin{exampleblock}{Ejemplo conceptual}
            ``Ejecutar estas instrucciones \textbf{10 veces}''
        \end{exampleblock}

        \column{0.5\textwidth}
        \textbf{Ciclo \texttt{while}}
        \begin{itemize}
            \item Se usa cuando \textbf{no se sabe} cu\'antas veces se repetir\'a.
            \item Repite mientras una \textbf{condici\'on sea verdadera}.
        \end{itemize}
        \vspace{0.3cm}
        \begin{exampleblock}{Ejemplo conceptual}
            ``Ejecutar estas instrucciones \textbf{mientras} el usuario no diga 'salir' ''
        \end{exampleblock}
    \end{columns}
\end{frame}

%================================================
\section{Sintaxis y estructura del \texttt{while}}
%================================================

\begin{frame}{Estructura del ciclo \texttt{while}}
    \begin{block}{Definici\'on}
        Un ciclo \texttt{while} ejecuta un bloque de c\'odigo \textbf{repetidamente}
        mientras que una \textbf{condici\'on l\'ogica} sea \texttt{True}.
        Cuando la condici\'on se vuelve \texttt{False}, el ciclo termina.
    \end{block}

    \vspace{0.3cm}
    \textbf{Sintaxis general en Python:}
    \[
    \underbrace{\texttt{while}}_{\text{palabra clave}}\;\;
    \underbrace{\texttt{condici\'on}}_{\text{expresi\'on booleana}}\texttt{:}
    \]
    \begin{itemize}
        \item La \textbf{condici\'on} se eval\'ua \emph{antes} de cada iteraci\'on.
        \item Si la condici\'on es \texttt{False} desde el principio, el cuerpo \textbf{nunca} se ejecuta.
        \item El cuerpo del ciclo debe \textbf{modificar la condici\'on} para evitar ciclos infinitos.
    \end{itemize}
\end{frame}

\begin{frame}{Diagrama de flujo del \texttt{while}}
    \begin{columns}
        \column{0.55\textwidth}
        \textbf{Flujo de ejecuci\'on:}
        \begin{enumerate}
            \item Se eval\'ua la \textbf{condici\'on}.
            \item Si es \texttt{True} $\rightarrow$ se ejecuta el \textbf{cuerpo}.
            \item Se regresa al paso 1.
            \item Si es \texttt{False} $\rightarrow$ el ciclo \textbf{termina}.
        \end{enumerate}

        \column{0.45\textwidth}
        \begin{alertblock}{Regla de oro}
            Siempre aseg\'urate de que la condici\'on eventualmente se vuelva \texttt{False},
            o usa \texttt{break} para salir del ciclo.
        \end{alertblock}
    \end{columns}
\end{frame}

%================================================
\section{\texttt{while} con condici\'on}
%================================================

\begin{frame}[fragile]{\texttt{while} con condici\'on expl\'icita}
    \begin{block}{Descripci\'on}
        La forma m\'as com\'un. El ciclo repite \textbf{mientras la condici\'on sea verdadera}.
        Normalmente se usa un \textbf{contador} o una variable que cambia en cada iteraci\'on.
    \end{block}

    \vspace{0.2cm}
    \textbf{Ejemplo 1 --- Contador ascendente:}
\begin{lstlisting}[style=PythonStyle]
contador = 1          # Inicializacion

while contador <= 5:  # Condicion
    print(f"Iteracion:{contador}")
    contador += 1     # Actualizacion (ESENCIAL)

print("Ciclo terminado.")
\end{lstlisting}
\end{frame}
\begin{frame}[fragile]{\texttt{while} con condici\'on expl\'icita}

    \begin{exampleblock}{Salida esperada}
        \texttt{Iteracion numero: 1} \quad
        \texttt{Iteracion numero: 2} \quad \ldots \quad
        \texttt{Iteracion numero: 5} \quad
        \texttt{Ciclo terminado.}
    \end{exampleblock}
\end{frame}

\begin{frame}[fragile]{Ejemplo 2 --- Suma acumulada}
    \textbf{Problema:} Sumar todos los n\'umeros del 1 al 100 usando \texttt{while}.

\begin{lstlisting}[style=PythonStyle]
n = 1
suma = 0

while n <= 100:
    suma += n   # suma = suma + n
    n += 1

print(f"La suma de 1 a 100 es: {suma}")
# Resultado: La suma de 1 a 100 es: 5050
\end{lstlisting}

    \begin{columns}
        \column{0.5\textwidth}
        \textbf{Variables clave:}
        \begin{itemize}
            \item \texttt{n}: contador (1 a 101)
            \item \texttt{suma}: acumulador
        \end{itemize}
        \column{0.5\textwidth}
        
    \end{columns}
\end{frame}

\begin{frame}[fragile]{Ejemplo 3 --- Validaci\'on de entrada}
    \textbf{Problema:} Pedir al usuario un n\'umero entre 1 y 10 hasta que lo ingrese correctamente.

\begin{lstlisting}[style=PythonStyle]
numero = int(input("Ingresa un numero entre 1 y 10: "))

while numero < 1 or numero > 10:
    print("Numero invalido. Intenta de nuevo.")
    numero = int(input("Ingresa un numero entre 1 y 10: "))

print(f"Numero aceptado: {numero}")
\end{lstlisting}

    \vspace{0.2cm}
    \begin{block}{Patr\'on: ciclo de validaci\'on}
        Este patr\'on es muy com\'un en aplicaciones reales. El ciclo act\'ua como
        un \textbf{guardabarrera}: solo deja pasar datos correctos.
    \end{block}
\end{frame}

\begin{frame}[fragile]{Ejemplo 4 --- Adivina el n\'umero}
    \textbf{Problema:} El programa elige el n\'umero secreto 7; el usuario adivina hasta acertar.

\begin{lstlisting}[style=PythonStyle]
secreto = 7
intento = int(input("Adivina el numero (1-10): "))
intentos = 1

while intento != secreto:
    if intento < secreto:
        print("Muy bajo. Intenta de nuevo.")
    else:
        print("Muy alto. Intenta de nuevo.")
    intento = int(input("Tu intento: "))
    intentos += 1

print(f"Correcto! Lo lograste en {intentos} intento(s).")
\end{lstlisting}
\end{frame}
\begin{frame}[fragile]{Ejemplo 4 --- Adivina el n\'umero}

    \begin{exampleblock}{Concepto reforzado}
        La condici\'on \texttt{intento != secreto} se eval\'ua cada vez.
        El ciclo termina exactamente cuando el usuario acierta.
    \end{exampleblock}
\end{frame}

%================================================
\section{\texttt{while True} y control con \texttt{break}}
%================================================

\begin{frame}{\texttt{while True}: ciclo infinito controlado}
    \begin{block}{Concepto}
        \texttt{while True} crea un ciclo \textbf{infinito intencionalmente}.
        La salida se controla con la instrucci\'on \texttt{break}, que interrumpe
        el ciclo desde cualquier punto dentro del cuerpo.
    \end{block}

    \vspace{0.3cm}
    \begin{columns}
        \column{0.5\textwidth}
        \textbf{?`Cu\'ando usarlo?}
        \begin{itemize}
            \item Men\'us interactivos.
            \item Servidores o daemons.
            \item Cuando la condici\'on de salida se calcula \emph{dentro} del cuerpo.
            \item Re-intentos hasta \'exito.
        \end{itemize}

        \column{0.5\textwidth}
        \begin{alertblock}{Diferencia clave}
            \textbf{\texttt{while condici\'on}:} la condici\'on se verifica \emph{al inicio}.\\[4pt]
            \textbf{\texttt{while True + break}:} la condici\'on de salida puede estar
            \emph{en cualquier lugar} del cuerpo.
        \end{alertblock}
    \end{columns}
\end{frame}

\begin{frame}[fragile]{\texttt{break} y \texttt{continue}}
    \begin{columns}
        \column{0.5\textwidth}
        \textbf{\texttt{break} --- salir del ciclo:}
\begin{lstlisting}[style=PythonStyle]
i = 0
while True:
    if i == 3:
        break  # Sale del ciclo
    print(i)
    i += 1
# Imprime: 0, 1, 2
\end{lstlisting}

        \column{0.5\textwidth}
        \textbf{\texttt{continue} --- saltar iteraci\'on:}
\begin{lstlisting}[style=PythonStyle]
i = 0
while i < 6:
    i += 1
    if i == 4:
        continue # Salta al siguiente
    print(i)
# Imprime: 1, 2, 3, 5, 6
\end{lstlisting}
    \end{columns}

    \vspace{0.3cm}
    \begin{block}{Resumen}
        \begin{itemize}
            \item \texttt{break}: termina el ciclo inmediatamente.
            \item \texttt{continue}: omite el resto del cuerpo en la iteraci\'on actual y pasa a la siguiente.
        \end{itemize}
    \end{block}
\end{frame}

\begin{frame}[fragile]{Ejemplo 5 --- Men\'u interactivo con \texttt{while True}}
    \textbf{Problema:} Men\'u de opciones que se repite hasta que el usuario elige salir.

\begin{lstlisting}[style=PythonStyle]
while True:
    print("\n--- MENU ---")
    print("1. Saludar")
    print("2. Mostrar la fecha")
    print("3. Salir")

    opcion = input("Elige una opcion: ")

    if opcion == "1":
        print("Hola, bienvenido!")
    elif opcion == "2":
        print("Hoy es lunes.")
    elif opcion == "3":
        print("Hasta luego.")
        break          # Salida controlada
    else:
        print("Opcion invalida.")
\end{lstlisting}
\end{frame}

\begin{frame}[fragile]{Ejemplo 6 --- Reintentos con \texttt{while True}}
    \textbf{Problema:} Solicitar contrase\~na hasta 3 intentos fallidos o \'exito.

\begin{lstlisting}[style=PythonStyle]
clave_correcta = "anahuac2025"
intentos = 0
MAX_INTENTOS = 3
while True:
    clave = input("Ingresa tu contrasena: ")
    intentos += 1

    if clave == clave_correcta:
        print("Acceso concedido.")
        break
    elif intentos >= MAX_INTENTOS:
        print("Cuenta bloqueada. Demasiados intentos.")
        break
    else:
        restantes = MAX_INTENTOS - intentos
        print(f"Clave incorrecta. Te quedan {restantes} intento(s).")
\end{lstlisting}
\end{frame}

\begin{frame}[fragile]{Ejemplo 6 --- Reintentos con \texttt{while True}}
    \begin{exampleblock}{Concepto reforzado}
        Dos condiciones de salida posibles dentro del cuerpo: \'exito o l\'imite de intentos.
    \end{exampleblock}
\end{frame}

%================================================
\section{Comparaci\'on y buenas pr\'acticas}
%================================================

\begin{frame}{Comparaci\'on: \texttt{while cond} vs \texttt{while True}}
    \begin{center}
    \begin{tabular}{lcc}
        \toprule
        \textbf{Aspecto} & \texttt{while condici\'on} & \texttt{while True + break} \\
        \midrule
        Condici\'on de salida & Al \textit{inicio} del ciclo & Dentro del cuerpo \\
        Legibilidad          & Alta (simple)               & Media (requiere \texttt{break}) \\
        Uso t\'ipico           & Contadores, rangos          & Men\'us, servidores, reintentos \\
        Riesgo ciclo infinito & Bajo (con cuidado)          & Alto sin \texttt{break} \\
        \bottomrule
    \end{tabular}
    \end{center}

    \vspace{0.3cm}
    \begin{alertblock}{Recomendaci\'on general}
        Prefiere \texttt{while condici\'on} cuando sea posible.
        Usa \texttt{while True} solo cuando la condici\'on de salida no puede evaluarse antes del cuerpo.
    \end{alertblock}
\end{frame}

\begin{frame}{Errores comunes y c\'omo evitarlos}
    \footnotesize
    \begin{enumerate}
        \item \textbf{Olvidar actualizar la variable de control}
        \begin{itemize}
            \item Causa un \textbf{ciclo infinito}.
            \item Soluci\'on: aseg\'urate de que \texttt{contador += 1} (o equivalente) est\'e en el cuerpo.
        \end{itemize}

        \vspace{0.2cm}
        \item \textbf{Condici\'on que nunca es \texttt{True} desde el inicio}
        \begin{itemize}
            \item El cuerpo \textbf{nunca} se ejecuta.
            \item Revisa los valores iniciales antes de entrar al ciclo.
        \end{itemize}

        \vspace{0.2cm}
        \item \textbf{Error de indentaci\'on}
        \begin{itemize}
            \item Python usa la \textbf{indentaci\'on} para delimitar el cuerpo del ciclo.
            \item Cada l\'inea del cuerpo debe tener exactamente 4 espacios de sangr\'ia.
        \end{itemize}

        \vspace{0.2cm}
        \item \textbf{Usar \texttt{while True} sin \texttt{break}}
        \begin{itemize}
            \item Programa colgado. Usa \texttt{Ctrl+C} para detenerlo.
        \end{itemize}
    \end{enumerate}
\end{frame}

%================================================
\section{Tareas}
%================================================

\begin{frame}{Tarea 1 --- Tabla de multiplicar interactiva}
    \begin{block}{Objetivo}
        Aplicar \texttt{while} con condici\'on para generar tablas de multiplicar y
        practicar la validaci\'on de entrada del usuario.
    \end{block}

    \vspace{0.2cm}
    \textbf{Instrucciones:}
    \begin{enumerate}
        \item Pide al usuario que ingrese un n\'umero entero entre 1 y 12.
              \textbf{Valida} la entrada: si est\'a fuera del rango, vuelve a pedirlo (usa \texttt{while}).
        \item Una vez validado, imprime la tabla de multiplicar completa del 1 al 12
              usando un ciclo \texttt{while}.
        \item Al terminar, pregunta: ``\textit{?`Deseas calcular otra tabla? (s/n)}''.
              Si el usuario responde \texttt{s}, repite el proceso. Si responde \texttt{n}, termina.
    \end{enumerate}
\end{frame}

\begin{frame}{Tarea 1 --- Tabla de multiplicar interactiva}

    \vspace{0.2cm}
    \begin{exampleblock}{Ejemplo de salida esperada (n\'umero = 7)}
        \texttt{7 x 1 = 7} \quad \texttt{7 x 2 = 14} \quad \ldots \quad \texttt{7 x 12 = 84}
    \end{exampleblock}

    \begin{alertblock}{Criterios de evaluaci\'on}
        Validaci\'on de entrada (30\%) \quad Ciclo de tabla (40\%) \quad Repetici\'on del programa (30\%)
    \end{alertblock}
\end{frame}

\begin{frame}{Tarea 2 --- Adivina el n\'umero (versi\'on completa)}
    \begin{block}{Objetivo}
        Combinar \texttt{while True} con \texttt{break} y el m\'odulo \texttt{random}
        para construir un juego interactivo completo.
    \end{block}

    \vspace{0.2cm}
    \textbf{Instrucciones:}
    \footnotesize
    \begin{enumerate}
        \item Genera un n\'umero secreto \textbf{aleatorio} entre 1 y 50 usando \texttt{random.randint(1, 50)}.
        \item Usa \texttt{while True} para pedir al usuario que adivine el n\'umero.
        \item En cada intento, indica si el n\'umero es ``demasiado alto'', ``demasiado bajo'' o ``correcto''.
        \item Cuando el usuario adivine, muestra el \textbf{n\'umero de intentos} utilizados y termina con \texttt{break}.
        \item \textbf{Mejora opcional (+20\% extra):} Agrega un l\'imite de 7 intentos.
              Si se agota, revela el n\'umero secreto y termina.
    \end{enumerate}
\end{frame}

\begin{frame}{Tarea 2 --- Adivina el n\'umero (versi\'on completa)}
    \begin{alertblock}{Criterios de evaluaci\'on}
        Generaci\'on aleatoria (20\%) \quad L\'ogica de pistas (40\%) \quad
        Conteo de intentos (20\%) \quad \texttt{break} correcto (20\%)
    \end{alertblock}
\end{frame}

%================================================
%   FINAL SLIDE
%================================================

\begin{frame}{Resumen}
    \begin{columns}
        \column{0.5\textwidth}
        \textbf{\texttt{while condici\'on}}
        \begin{itemize}
            \item Condici\'on al inicio.
            \item Ideal para contadores y rangos.
            \item Requiere actualizar la variable de control.
        \end{itemize}
        \vspace{0.3cm}
        \textbf{\texttt{while True + break}}
        \begin{itemize}
            \item Salida desde dentro del cuerpo.
            \item Ideal para men\'us y reintentos.
            \item Siempre debe haber un \texttt{break}.
        \end{itemize}

        \column{0.5\textwidth}
        \begin{exampleblock}{Comandos adicionales}
            \begin{itemize}
                \item \texttt{break}: salir del ciclo.
                \item \texttt{continue}: saltar iteraci\'on actual.
            \end{itemize}
        \end{exampleblock}
        \vspace{0.3cm}
        \begin{alertblock}{Recuerda}
            Todo ciclo \texttt{while} debe tener
            una condici\'on de salida garantizada.
        \end{alertblock}
    \end{columns}
\end{frame}

\end{document}